\documentclass[a4paper]{umbruch-position}
\usepackage[utf8]{inputenc}
\usepackage[T1]{fontenc}
\usepackage[ngerman]{babel}
\usepackage{tikz}
\usepackage{amssymb}
\usepackage{eurosym}
\usepackage{graphicx}
\usepackage{xcolor}
\definecolor{umbruch-dark}{HTML}{2A2A2A}
\definecolor{umbruch-teal}{HTML}{008080}
\definecolor{umbruch-red}{HTML}{B22222}
\usetikzlibrary{shapes.geometric, arrows.meta, positioning, fit, calc}

\positionsnummer{Basiert auf der Studie ``Systemtheoretische Aufarbeitung der Regressangst'' (ST-001)}
\versionsnummer{0.2}
\title{Der Architekturplan gegen den Regress}
\untertitel{Ein strategisches Protokoll f\"ur niedergelassene \"Arzte:\\Pr\"avention, Firewall-Aufbau und juristische Reaktion}

\begin{document}

\maketitle

\begin{infobox}
\textbf{WICHTIGER HINWEIS:}\\
\textit{Dieser Leitfaden dient der allgemeinen Information und stellt keine Rechtsberatung dar. F\"ur individuelle rechtliche Beratung wenden Sie sich an einen Fachanwalt f\"ur Sozialrecht (UPD, 0800 011 77 22).}
\end{infobox}

\section{Ruhe bewahren --- Der Kontext}

Sie haben einen Pr\"ufungsbescheid erhalten. Atmen Sie durch.
Die Wahrscheinlichkeit daf\"ur lag bei unter 1\,\% (beide Pr\"ufverfahren zusammen; das statistische Screening dominiert diese Zahl), doch nun ist es passiert. Wichtig ist: Viele Bescheide werden im Rahmen des Widerspruchsverfahrens reduziert oder komplett aufgehoben!

Sie sind nicht allein. Die Androhung kann existenzgef\"ahrdend wirken, oft reduziert sich der Schadenswert jedoch bei solider Dokumentation erheblich. Handeln Sie fristgerecht.

\section{Was dieser Bescheid sagt~--- und was nicht}

Bevor Sie handeln: Drei Gleichsetzungen, die im Verfahren gef\"ahrlich sind:

\begin{center}
{\small
\begin{tabular}{p{5cm} p{7cm}}
\toprule
\textbf{Nicht gleichzusetzen} & \textbf{Warum} \\
\midrule
Verdacht unbegr\"undet $\neq$ medizinisch korrekt & Widerspruchserfolg basiert auf Verfahrensfehlern, nicht auf medizinischer Neubewertung \\[4pt]
Widerspruch erfolgreich $\neq$ medizinisch zertifiziert & Formale Entlastung~--- das Widerspruchsgremium pr\"uft Rechtm\"a\ss{}igkeit, keine originäre medizinische Qualit\"at \\[4pt]
Klage gewonnen $\neq$ Pr\"azedenzschutz & Entscheidung gilt im Einzelfall; keine Schutzwirkung f\"ur sp\"atere Verfahren \\
\bottomrule
\end{tabular}
}
\end{center}

Und: Je nachdem welches Verfahren Sie betrifft, unterscheiden sich Ihre Chancen erheblich:

{\small
\begin{tabularx}{\textwidth}{p{3.5cm} X X}
\toprule
 & \textbf{Auff\"alligkeitspr\"ufung} & \textbf{Einzelfallpr\"ufung} \\
\midrule
Beratung vor Regress & Vorgeschrieben & Ausgeschlossen \\
Betroffenheit & Teilweise & 100\,\% \\
Neg.\,Ausgang & Niedrig wahrscheinlich & Hoch wahrscheinlich\footnote{Gilt prim\"ar f\"ur V2b (Formfehler). Bei V2a (Unwirtschaftlichkeit) ist eine medizinische Verteidigung m\"oglich.} \\
Pos.\,Ausgang & Hoch wahrscheinlich & Mittel bis niedrig \\
Schaden\-h\"ohe & Niedrig bis moderat & Moderat bis hoch \\
Existenz\-bedrohend & Sehr gering & Niedrig, aber real \\
\bottomrule
\end{tabularx}
}

\textbf{Das ist kein akademischer Punkt.} Wer versteht, dass das System \textit{Verfahrensfehler}, nicht \textit{medizinische Angemessenheit} pr\"uft, kann seinen Widerspruch gezielt auf das Richtige ausrichten~--- dies gilt uneingeschr\"ankt f\"ur V2b (Formfehler). Bei V1 und V2a k\"onnen medizinische Gr\"unde den Ausgang entscheidend beeinflussen.

\section{Entscheidungsbaum: \glqq{}Pr\"ufungsbescheid erhalten\grqq{}}

\begin{infobox}
\textit{Entscheidungshilfe, keine Rechtsberatung. Fristen und Rechtsgrundlagen k\"onnen sich \"andern.}
\end{infobox}

\begin{center}
\tikzset{
    action/.style={rectangle, draw=umbruch-dark, fill=umbruch-teal!15, rounded corners, minimum height=1cm, align=center, text width=4cm, inner sep=0.2cm},
    decision/.style={diamond, aspect=2.0, draw=umbruch-dark, fill=white, align=center, text width=3cm, inner sep=0cm},
    warning/.style={rectangle, draw=umbruch-red, fill=umbruch-red!15, rounded corners, minimum height=1cm, align=center, text width=4cm, inner sep=0.2cm},
    arrow/.style={-{Stealth[scale=1.2]}, thick, umbruch-dark}
}
\begin{tikzpicture}[
    scale=0.8, transform shape,
    node distance=1.5cm and 1cm,
    font=\sffamily\small
]
\node[warning] (bescheid) {Bescheid erhalten!\\Frist: 1 Mnt. (§ 84 SGG)};
\node[decision, below=1cm of bescheid] (art) {Auff\"alligkeits- oder\\Einzelfallpr\"ufung?};
\node[decision, left=1cm of art] (auffaellig) {Praxisbeson-\\derheiten da?};
\node[decision, right=1cm of art] (einzelfall) {Leitlinien-\\gerecht?};
\node[action, left=1cm of auffaellig] (auff_ja) {Widerspruch mit Begr\"undung};
\node[action, below=1.5cm of auffaellig] (auff_nein) {Beratung anfor-\\dern, Stellungnahme};
\node[action, right=1cm of einzelfall] (einz_ja) {Widerspruch mit\\Leitlinie \& Akte};
\node[action, below=1.5cm of einzelfall] (einz_nein) {Teilanerkennung/\\Beratung beantragen};
\node[action, below=4cm of art] (widerspruch) {Widerspruch frist-gerecht einlegen!};
\node[decision, below=1cm of widerspruch] (ausschuss) {Beschwerde-\\ausschuss};
\node[action, left=1cm of ausschuss] (reduziert) {Regress\\aufgehoben};
\node[warning, right=1cm of ausschuss] (gericht) {Klage SG!\\Frist: 1 Mnt (§ 87 SGG)};
\draw[arrow] (bescheid) -- (art);
\draw[arrow] (art) -- node[above] {Auff\"allig} (auffaellig);
\draw[arrow] (art) -- node[above] {Einzelfall} (einzelfall);
\draw[arrow] (auffaellig) -- node[above] {JA} (auff_ja);
\draw[arrow] (auffaellig) -- node[left] {NEIN} (auff_nein);
\draw[arrow] (einzelfall) -- node[above] {JA} (einz_ja);
\draw[arrow] (einzelfall) -- node[left] {NEIN} (einz_nein);
\draw[arrow] (auff_ja) |- (widerspruch);
\draw[arrow] (auff_nein) |- (widerspruch);
\draw[arrow] (einz_ja) |- (widerspruch);
\draw[arrow] (einz_nein) |- (widerspruch);
\draw[arrow] (widerspruch) -- (ausschuss);
\draw[arrow] (ausschuss) -- node[above] {Erfolg} (reduziert);
\draw[arrow] (ausschuss) -- node[above] {Abgelehnt} (gericht);
\end{tikzpicture}
\end{center}

\newpage
\section{Haben Sie pr\"aventive Ma\ss{}nahmen getroffen? Diese helfen jetzt}

Falls Sie vor der strittigen Verordnung eine \textbf{KV-Verordnungsberatung} in Anspruch genommen haben oder eine \textbf{Stellungnahme der Krankenkasse zur Verordnungsf\"ahigkeit} (BSG B\,6\,KA\,27/12\,R) eingeholt haben, bringen Sie diese Dokumente jetzt zwingend in Ihren Widerspruch ein:

\begin{itemize}
    \item \textbf{KV-Beratungsprotokoll:} Belegt, dass Sie die Verordnung nicht eigenmächtig, sondern auf Basis fachlicher KV-Beratung getroffen haben. Das st\"arkt Ihre Position gegen\"uber der Pr\"ufungsstelle erheblich.
    \item \textbf{KK-Stellungnahme zur Verordnungsf\"ahigkeit:} Wenn die Kasse die Verordnungsf\"ahigkeit vorab best\"atigt hat, kann sie als Antragstellerin eines Einzelfallregresses nicht widerspruchsfrei handeln. Dieses Argument geh\"ort in den Widerspruch.
\end{itemize}

Falls Sie diese Ma\ss{}nahmen nicht getroffen haben: F\"ur k\"unftige Verordnungen k\"onnen Sie beide pr\"aventiven Wege nutzen (siehe LF-001).

\section{Fristen und Formalia}

Wenn ein Pr\"ufungsbescheid vorliegt, gelten harte gesetzliche Fristen.

\subsection*{Ihre Reaktionsfristen}
\begin{itemize}
    \item \textbf{Widerspruchsfrist:} Sie betr\"agt genau \textbf{1 Monat} ab Feststellung bzw. Zustellung (\S\,84 SGG)! Reichen Sie den Widerspruch vorab ``dem Grunde nach'' (zur Fristwahrung) per Einschreiben ein, eine Begr\"undung kann nachgereicht werden.
    \item \textbf{Fehlende Rechtsbehelfsbelehrung:} Ist dem Bescheid keine oder eine fehlerhafte Belehrung beigef\"ugt, verl\"angert sich die Widerspruchsfrist auf \textbf{1 Jahr} (\S\,66 Abs.\,2 SGG).
\end{itemize}

\subsection*{Wie weit zur\"uck d\"urfen die Sie pr\"ufen? (Ausschlussfristen)}

Entscheidend: Es handelt sich um \textbf{Ausschlussfristen} (das Recht erlischt automatisch nach Ablauf --- Sie m\"ussen nichts einwenden). Das ist die f\"ur Sie g\"unstigere Variante gegen\"uber einer Verj\"ahrung.

{\small
\begin{tabularx}{\textwidth}{l X X}
\toprule
 & \textbf{V1 Auff\"alligkeitspr\"ufung} & \textbf{V2 Einzelfallpr\"ufung} \\
\midrule
\multicolumn{3}{l}{\textit{Aktuelles Recht (seit TSVG 2019 / GVWG 2021)}} \\
\addlinespace[2pt]
Frist & \textbf{2 Jahre} ab Jahresende & \textbf{2,5 Jahre} (18 Mon. Antragsfrist + 12 Mon. Festsetzung) \\
\addlinespace
Rechtsgrundlage & \S\,106 Abs.\,3 S.\,3 SGB\,V & \S\,106 Abs.\,3 S.\,4--5 SGB\,V \\
\midrule
\multicolumn{3}{l}{\textit{Altes Recht (vor 11.05.2019 --- gilt f\"ur Altf\"alle!)}} \\
\addlinespace[2pt]
Frist & \textbf{4 Jahre} & \textbf{4 Jahre} \\
\addlinespace
Rechtsgrundlage & BSG 6\,RKa\,40/94 & BSG B\,6\,KA\,13/18\,R \\
\bottomrule
\end{tabularx}
}

\begin{infobox}
\textbf{Praxishinweis:} Pr\"ufen Sie als erstes, ob der Bescheid die Ausschlussfrist einhält! Betrifft der Bescheid Verordnungen, die \textbf{mehr als 2,5 Jahre} zur\"uckliegen (bzw. 4 Jahre bei Altf\"allen vor 2019), ist er m\"oglicherweise verfristet. Wenn ja: Im Widerspruch \textbf{sofort} r\"ugen. --- \textbf{Aber Achtung:} Wurde das Verfahren rechtzeitig eingeleitet, gibt es f\"ur den anschlie\ss{}enden Instanzenzug (Beschwerdeausschuss, Sozialgericht, LSG, BSG) keine zeitliche Obergrenze. Verfahren mit 6--10 Jahren Gesamtdauer sind dokumentiert.
\end{infobox}

\section{Widerspruch formulieren --- Textbausteine}

\begin{infobox}
\textbf{Baustein 1: Einlegung des Widerspruchs (dem Grunde nach)}\\
\textit{\glqq{}Gegen Ihren Bescheid vom [Datum], Az. [Nummer], lege ich hiermit form- und fristgerecht Widerspruch ein (§ 84 SGG). Die ausf\"uhrliche Begr\"undung unter Nachreichung der entlastenden Patientenakten erfolgt in k\"urze.\grqq{}}
\end{infobox}

\begin{infobox}
\textbf{Baustein 2: Beratungsantrag (\S\,106b Abs. 2 SGB V) --- NUR bei Auff\"alligkeitspr\"ufung (V1)!}\\
\textit{\glqq{}Da es sich um Auff\"alligkeiten im Rahmen einer erstmaligen Pr\"ufung handelt, beantrage ich hiermit die Durchf\"uhrung einer eingehenden KV-Beratung anstelle einer sofortigen Honorark\"urzung gem\"a\ss{} \S\,106b Abs.\,2 SGB\,V.\grqq{}}\\[4pt]
\textbf{Achtung:} Dieser Baustein gilt \textbf{nicht} bei Einzelfallpr\"ufungen (V2). Dort ist Beratung vor Regress gesetzlich ausgeschlossen (\S\,106b Abs.\,2 S.\,4). Bei V2: direkt Widerspruch mit Dokumentation einlegen.
\end{infobox}

\section{Wichtig: Formfehler sind im Regressverfahren nicht heilbar}

Eine zentrale Erkenntnis aus der j\"ungsten Rechtsprechung: \textbf{Formfehler bei Verordnungen sind im Regressverfahren nach st\"andiger BSG-Rechtsprechung nicht nachtr\"aglich heilbar.} Im Unterschied zum Steuerrecht (\S\S\,129, 153 AO), zum Verwaltungsrecht (\S\,45 VwVfG) und zum Baurecht (\S\S\,214--215 BauGB) sieht das Sozialgesetzbuch f\"ur das Regressverfahren keine Heilung vor.

\begin{infobox}
\textbf{Wichtige Unterscheidung:} \textit{Vor} Einleitung eines Pr\"ufverfahrens k\"onnen Formfehler noch korrigiert werden (z.\,B. fehlende Unterschrift nachtragen, Angaben erg\"anzen). \textit{Nach} Einleitung greift der sogenannte \textbf{normative Schadensbegriff} (BSG B\,6\,KA\,5/09\,R): Der Schaden der Kasse gilt ab dem Moment als eingetreten, in dem sie die formal fehlerhafte Verordnung bezahlt hat --- unabh\"angig davon, ob das Medikament medizinisch notwendig war. Ab diesem Zeitpunkt ist Heilung ausgeschlossen. Das Verfahren pr\"uft nur noch: Lag der Formfehler zum Zeitpunkt der Verordnung vor? Ja oder Nein.
\end{infobox}

Das BSG hat dies am 27.08.2025 (B\,6\,KA\,9/24\,R) bei einem Stempel-Regress \"uber ca. 491.000\,EUR best\"atigt: Die pers\"onliche Unterschrift sei kein blo\ss{}er Formalakt, sondern diene dem ``Schutz des Lebens''. Der Arzt argumentierte, alle Verordnungen seien medizinisch indiziert gewesen. \textbf{Das Gericht erkl\"arte dieses Argument f\"ur irrelevant.} Die KBV nannte das Urteil ``unverh\"altnism\"a\ss{}ig und \"uberzogen''. Der betroffene Arzt k\"undigte eine Verfassungsbeschwerde an.

\textbf{Was das f\"ur Sie bedeutet:} Wenn ein Formfehler festgestellt wird und das Verfahren eingeleitet ist, steht das Ergebnis im Wesentlichen fest. Deshalb lassen sich gesch\"atzt 70--80\,\% der betroffenen \"Arzte auf einen \textbf{Vergleich} ein (typisch 30--50\,\% des Regressbetrags), statt das Verfahren bis zum Sozialgericht zu f\"uhren. Wer zum Gericht geht in der Erwartung, das Gericht w\"urde die medizinische Korrektheit der Behandlung w\"urdigen, wird entt\"auscht: \textbf{Bei V2b (Formfehler) pr\"uft das Gericht die Form, nicht die Medizin. Bei V1 und V2a kann medizinische Argumentation vor Gericht erheblichen Einfluss haben.}

\textbf{Konsequenz f\"ur den Widerspruch:} Bei einem Formalregress hat die Kasse f\"ur eine medizinisch korrekte Verordnung bezahlt. H\"atte der Arzt formfehlerfrei verordnet, h\"atte die Kasse \textit{denselben Betrag} zahlen m\"ussen. Der tats\"achliche Mehrschaden ist Null --- dennoch droht voller Regress. Das BSG hat die Differenzkostenregelung (\S\,106b Abs.\,2a) 2024 f\"ur unwirtschaftliche Verordnungen anerkannt (B\,6\,KA\,5/23\,R, B\,6\,KA\,10/23\,R), f\"ur \textbf{unzul\"assige} Verordnungen (= Formfehler) aber ausdr\"ucklich ausgeschlossen. Das Eigentums- und Verh\"altnism\"a\ss{}igkeitsargument bleibt Teil der Verteidigungslinie.

Diese rechtsdogmatische Schieflage ist Gegenstand der angek\"undigten Verfassungsbeschwerde.

\section{Wenn der Widerspruch scheitert: Sozialgericht}

Wird der Widerspruch abgewiesen, bleibt der Klageweg vor dem Sozialgericht (Klagefrist \textbf{1 Monat}, § 87 SGG).
Erfolgsaussichten sind oft gut bei formellen Verfahrensfehlern und stichhaltigen Praxisbesonderheiten (wie dokumentierten schweren Patientenf\"allen).
F\"ur Vertrags\"arzte fallen beim Sozialgericht Gerichtsgeb\"uhren nach GKG an (\S\,197a SGG --- anders als f\"ur Versicherte, die gerichtskostenfrei klagen). Die Kosten sind kalkulierbar, zus\"atzlich m\"ussen noch eigene Anwaltskosten von ca. 1.500--5.000 EUR getragen werden, sofern keine Regress-Rechtsschutzversicherung besteht.

\section{Vier F\"alle --- Wie es anderen erging}

\begin{itemize}
    \item \textbf{Fall 1: Physiotherapeut} -- 65.000 EUR Forderung wegen Durchschnitts\"uberschreitung. Bescheid wurde reduziert, weil schwere neurologische Erkrankungen des Klientels (Praxisbesonderheit) best\"atigt wurden.
    \item \textbf{Fall 2: Heilmittelregress Schwergewichte} -- Existenzbedrohend durch sechsstellige Forderung. Das absurde Hindernis: Fehlende Versichertennummern im Bescheid, obwohl diese vorlagen. Widerspruch läuft!
    \item \textbf{Fall 3: Cannabis-Vaporizer} -- Kasse genehmigt teuren Vaporizer, regressiert dann aber die Cannabis-Bl\"uten. Hier greift das Vertrauensparadox. 
    \item \textbf{Fall 4: Unzul\"assige Delegation} -- Die Regressstelle forderte Gelder zur\"uck, da delegationsf\"ahige Hausbesuche f\"alschlich als arzt-durchgef\"uhrt (EBM-Ziffern) kodiert wurden.
\end{itemize}

\section{Langfristige Lehren: Struktur statt Angst}

Dieser Regress ist ein Datenpunkt~--- kein medizinisches Urteil. Drei Erkenntnisse f\"ur die Zeit danach:

\begin{enumerate}
    \item \textbf{Regress $\neq$ medizinisches Fehlverhalten.} Das System pr\"uft \"uberwiegend formale Angreifbarkeit, nicht medizinische Qualit\"at (${\sim}$35\,\% der Pr\"ufungen codieren allerdings medizinische Evidenz als Formregeln --- Biosimilar-Pflicht, AM-RL). Wer formal angreifbar war, verliert~--- auch wenn die Verordnung leitliniengerecht war.
    \item \textbf{Pr\"avention beginnt vor dem Antrag.} Einzelfallverfahren entstehen, weil Kassen einen begr\"undeten Antrag stellen k\"onnen. Wer Praxisbesonderheiten fr\"uhzeitig dokumentiert und die KV-Beratung nutzt, reduziert die Angriffsfl\"ache strukturell. Bei \textit{unklarer Erstattungsf\"ahigkeit} (Off-Label, neue Wirkstoffe) kann eine KK-Stellungnahme (BSG B\,6\,KA\,27/12\,R) Vertrauensschutz schaffen --- bei klar verordnungsf\"ahigen Standardmedikamenten ist sie inhaltsleer.
    \item \textbf{Psychologische Realit\"at anerkennen.} Einzelfallverfahren erzeugen Druck, Wiederholungsangst und das Gef\"uhl des Kontrollverlusts. Das ist rational~--- V2 hat keine eingebauten Schutzstufen~--- anders als V1, das ein Stufenmodell mit Beratung, Karenzzeit und Regressdeckel (25.000\,EUR) vorsieht. Wer das versteht, kann gezielt gegensteuern statt defensiv zu verordnen.
\end{enumerate}

\newpage
\section{Checkliste: Sofort nach Pr\"ufungsbescheid}

\begin{center}
\textbf{Was ist in den ersten 48 Stunden zu tun?}
\end{center}

\begin{itemize}
    \item[$\square$] \textbf{Frist notieren} (Exakt 1 Monat nach Zustellung, § 84 SGG).
    \item[$\square$] \textbf{Bescheid vollst\"andig lesen} (Inklusive aller Anlagen und der Rechtsbehelfsbelehrung).
    \item[$\square$] \textbf{Widerspruch ``dem Grunde nach'' einlegen} (Zur reinen Fristwahrung per Einschreiben R\"uckschein).
    \item[$\square$] \textbf{KV-Beratung kontaktieren} (Beratung nach \S\,106b Abs.\,2 voll aussch\"opfen --- \textbf{Achtung: gilt nur bei Auff\"alligkeitspr\"ufung!} Bei Einzelfallpr\"ufung besteht kein gesetzlicher Beratungsanspruch; hier direkt Widerspruch mit Dokumentation einlegen).
    \item[$\square$] \textbf{Rechtsschutzversicherung / Regressschutz} (Z.\,B. Therapiefreiheit e.V. oder Virchowbund infomieren).
    \item[$\square$] \textbf{Beweissicherung der Patientenakten} (Exakte Exzerpte vorbereiten, die Leitlinienkonformit\"at oder Praxisbesonderheiten sichern).
\end{itemize}

\vspace{1cm}

\begin{infobox}
\textbf{Gemeinsamer Rahmen (gilt f\"ur alle drei Leitf\"aden):}\\
Die ``unter 1\,\%''-Regressquote z\"ahlt nur V1-Regresse, nur rechtskr\"aftig. Bundesweit: ca.\,47.000 V2-Verfahren (BMG 2022). In Bayern: 13.332 Regresse bei $\sim$24.000 \"Arzten (SpiFa 2024). Der Faktor zwischen offizieller Quote und tats\"achlicher Erfahrung: ca.\,500$\times$. Beide Zahlen stimmen --- sie messen verschiedene Dinge. \textbf{Kleines Risiko pro Fall summiert sich:} 3.000 Verordnungen/Quartal = 120.000 Angriffspunkte in 10 Jahren. Das erkl\"art Ribbat: 72\,\% eigene Regresserfahrung. Sicherheit entsteht durch Struktur --- nicht durch Zur\"uckhaltung.
\end{infobox}

\vfill
\begin{center}
\footnotesize{\textit{Dieser Leitfaden dient der allgemeinen Information und stellt keine Rechtsberatung dar. F\"ur individuelle rechtliche Beratung wenden Sie sich an einen Fachanwalt f\"ur Sozialrecht oder die Unabh\"angige Patientenberatung (UPD, 0800 011 77 22).}}
\end{center}


\end{document}
